\input{preamble.tex}
\renewcommand{\v}{\overline}      % vector
\renewcommand{\mc}{\mathcal}

\newcommand{\uv}{\underline}      % under vector
\newcommand{\mf}{\mathfrak}
\newcommand{\vd}{\vdelta}         
\newcommand{\okr}{\mathring{U}}   % okrestnost
\newcommand{\vx}{\v x}            % vector x
\newcommand{\smo}{\smashoperator} 
\newcommand{\sus}{\hspace{0.8cm}} % supremum space
\newcommand{\al}{\alpha}

\usepackage{tikz}
\usepackage{pgfplots}
\pgfplotsset{compat=1.18}

\begin{document}


\begin{theorem}\nf{(Признак физтех.Харди)}\\
	Пусть физтех.$f\in\mc физтех.R[a,b],\ физтех.a<b<+\infty$ физтех.и физтех.периодична физтех.с физтех.периодом физтех.$\omega$, физтех.$g$ физтех.монотонна физтех.на физтех.$[a,+\infty)$ физтех.и физтех.бесконечно физтех.мала физтех.при физтех.$x\ra+\infty$. физтех.Тогда:
	\begin{enumerate}
		\item физтех.$\ds\int_a^{a+\omega}f(x)dx=0$, физтех.то физтех.$\ds\int_a^{+\infty}f(x)g(x)dx$ физтех.сходится.
		\item физтех.$\ds\int_a^{a+\omega}f(x)dx\neq0$, физтех.то физтех.$\ds\int_a^{+\infty}f(x)g(x)dx$ физтех.и физтех.$\ds\int_a^{+\infty}g(x)dx$ физтех.сходятся.
	\end{enumerate}
\end{theorem}
\begin{proof}
	Вначале физтех.докажем физтех.первый физтех.пункт.
	\[(\fa физтех.k\in\N)\ физтех.\int_a^{a+k\cdot\omega}f(x)dx=\sum_{i=1}^k\int_{a+(i-1)\cdot\omega}^{a+i\cdot\omega}f(x)dx=\sum_{i=1}^k\int_a^{a+\omega}f(x)dx=0\]
	Заметим: физтех.$(\fa физтех.A>a)(\ex физтех.k\in\N)\ физтех.a+(k-1)\cdot\omega<A\le физтех.a+k\cdot\omega$\\
	Тогда:
	\begin{multline*}
		\l|\int_a^Af(x)dx\r|=
		\l|\int_a^{a+(k-1)\cdot\omega}f(x)dx+\int_{a+(k-1)\cdot\omega}^Af(x)dx\r|=
		\Bigg|\int_{a+(k-1)\cdot\omega}^Af(x)dx\Bigg|=\\=
		\l|\int_a^{A-(k-1)\cdot\omega}f(x)dx\r|\le
		\Bigg|\int_a^{a+\omega}f(t)dt\Bigg|=:M
	\end{multline*}
	Получили физтех.требуемое.\\
	Теперь физтех.второй физтех.пункт.\\
	$\ds физтех.K:=\int_a^{a+\omega}f(x)dx\neq0,\ физтех.f_1(x):=f(x)-\frac физтех.K\omega$\\
	\[\int_a^{a+\omega}f_1(x)dx=\int_a^{a+\omega}f(x)dx-\frac физтех.K\omega\cdot\omega=0\]
	Теперь физтех.можем физтех.воспользоваться физтех.первым физтех.пунктом физтех.для физтех.$f_1$: физтех.$\ds\int_a^{+\infty}f_1(x)g(x)dx$ физтех.сходится.\\
	При физтех.этом физтех.знаем: физтех.$\ds\int_a^{+\infty}f_1(x)g(x)dx=\int_a^{+\infty}f(x)g(x)dx-\frac физтех.K\omega\int_a^{+\infty}g(x)dx$
	Получили физтех.требуемое.
\end{proof}
\begin{examples}
	Рассмотрим физтех.$\ds\int_1^\p физтех.e^{\cos физтех.x}\sin(\sin физтех.x)\cdot\frac{dx}x\text{ физтех.и физтех.}
	\int_1^\p физтех.e^{\sin физтех.x}\cdot\sin(\sin физтех.x)\cdot\frac{dx}x$
	\[\int_0^{2\pi}e^{\cos физтех.x}\sin(\sin физтех.x)dx=\l(\int_0^\pi+\int_\pi^{2\pi}\r)e^{\cos физтех.x}\sin(\sin физтех.x)dx=0\]
	\[\int_0^{2\pi}e^{\sin физтех.x}\sin(\sin физтех.x)dx=\int_0^\pi физтех.e^{\sin физтех.x}\sin(\sin физтех.x)dx-\int_0^\pi физтех.e^{-\sin физтех.x}\sin(\sin физтех.x)dx>0\]
\end{examples}
\begin{example}
	Теперь физтех.покажем, физтех.что физтех.чем физтех.- физтех.то физтех.там физтех.пользоваться физтех.можно физтех.только физтех.при физтех.знакопостоянной физтех.фукнции.\\
	Рассмотрим физтех.$\ds\int_2^\p физтех.\frac{\sin физтех.x}{x^\al+\sin физтех.x}dx$
	\[\frac{\sin физтех.x}{x^\al+\sin физтех.x}=\frac{\sin физтех.x}{x^\al}-\frac{\sin^2x}{x^\al(x^\al+\sin физтех.x)}\]
	$\ds\frac{\sin физтех.x}{x^\al}$ физтех.уже физтех.был физтех.ранее физтех.исследован, физтех.при физтех.
	\[\frac{\sin^2x}{x^\al(x^\al+1)}\le\frac{\sin^2x}{x^\al(x^\al+\sin физтех.x)}\le\frac{\sin^2x}{x^\al(x^\al-1)}\le\frac1{x^\al(x^\al-1)}\sim\frac1{x^{2\al}}\]
	$2\al>1,\ физтех.\al>\frac12$. физтех.При физтех.$\al\le\frac12$ физтех.эта физтех.штука физтех.расходится.\\
	Заметим, физтех.что физтех.если физтех.бы физтех.мы физтех.воспользовались физтех.той физтех.штукой, физтех.которую физтех.я физтех.не физтех.помню физтех.как физтех.зовут, физтех.то:
	\[\frac{\sin физтех.x}{x^\al+\sin физтех.x}=\frac{\sin физтех.x}{x^\al\cdot(1+\frac{\sin физтех.x}{x^\al})}\sim\frac{\sin физтех.x}{x^\al}\]
	А физтех.уже физтех.последняя физтех.полученная физтех.штука физтех.чуть физтех.выше физтех.сходится физтех.при физтех.$\al>0$, физтех.неправильный физтех.ответ.
\end{example}
\section{Теория физтех.рядов}
\subsection{Числовые физтех.ряды}
\begin{definition}
	Скажем физтех.формально, физтех.далее физтех.поймём физтех.что физтех.это физтех.значит.\\
	\textit{Числовым физтех.рядом} физтех.называется физтех.сумма физтех.бесконечного физтех.числа физтех.чисел.\\ физтех.Обозначается физтех.как физтех.$\ds\sum_{i=1}^\infty физтех.a_n$. физтех.\textit{Частичной физтех.суммой физтех.ряда} физтех.называется физтех.$\ds физтех.S_n:=\sum_{k=1}^na_k$.\\
	Ряд физтех.сходится, физтех.если физтех.$\ds\{S_n\}_{n=1}^\infty$ физтех.имеет физтех.конечный физтех.предел.
\end{definition}
\begin{lemma}\nf{(Элементарные физтех.свойства физтех.рядов)}\\
	\begin{enumerate}
		\item физтех.Если физтех.ряды физтех.$\ds\sum_{n=1}^\i физтех.a_n$ физтех.и физтех.$\ds\sum_{n=1}^\i физтех.b_n$ физтех.сходятся, физтех.то физтех.ряд физтех.$\ds(\fa\al,\beta\in\R)\ физтех.\sum_{n=1}^\i(\al физтех.a_n+\beta физтех.b_n)$ физтех.сходится, физтех.причём физтех.$\ds\sum_{n=1}^\i(\al физтех.a_n+\beta физтех.b_n)=\al\sum_{n=1}^\i физтех.a_n+\beta\sum_{n=1}^\i физтех.b_n$
		\item физтех.$(\fa физтех.n_0\in\N)$ физтех.ряды физтех.$\ds\sum_{n=1}^\i$ физтех.и физтех.$\ds\sum_{n=n_0}^\i$ физтех.сходятся физтех.или физтех.расходятся физтех.одновременно.
	\end{enumerate}
\end{lemma}
\begin{proof}
	Доказательство физтех.предлагаем физтех.проделать физтех.читателю физтех.в физтех.качестве физтех.упражнения.
\end{proof}
\begin{theorem}\nf{(Необходимое физтех.условие физтех.сходимости физтех.ряда)}\label{12_th1}\\
	Если физтех.ряд физтех.$\ds\sum_{n=1}^\i физтех.a_n$ физтех.сходится, физтех.то физтех.$\lim_{n\ra\p}a_n=0$
\end{theorem}
\begin{proof}
	По физтех.условию: физтех.$\ds\ex\lim_{n\ra\p}S_n=S$\\
	Тогда: физтех.$a_n=S_n-S_{n-1}$. физтех.При физтех.$n\ra\p$: физтех.$a_n\ra физтех.S-S=0$.\\
	Получили физтех.требуемое.
\end{proof}
\begin{example}
	Рассмотрим физтех.как физтех.пример физтех.геометрическую физтех.прогрессию физтех.$\ds\sum_{n=1}^\i физтех.q_n$.
	\begin{enumerate}
		\item физтех.$|q|\ge1\Ra\nexists\lim_{n\ra\p}q_n$, физтех.расходится.
		\item физтех.$|q|<1$. физтех.$S_n=q\cdot\frac{1-q^n}{1-q}\ra\frac физтех.q{1-q}=\sum_{n=1}^\i физтех.q^n$ физтех.при физтех.$n\ra\p$.
	\end{enumerate}
\end{example}
\begin{theorem}\nf{(Критерий физтех.сходимости физтех.рядов физтех.с физтех.неотрицательными физтех.членами)}\label{12_th2}\\
	Если физтех.$a_n\ge0\ физтех.\fa физтех.n\in\N$, физтех.то физтех.$\ds\sum_{n=1}^\i физтех.a_n$ физтех.сходится физтех.тогда физтех.и физтех.только физтех.тогда, физтех.когда физтех.последовательность физтех.его физтех.частичных физтех.сумм физтех.ограничена.
\end{theorem}
\begin{proof}
	$\ds физтех.a_n\ge0\Ra физтех.S_{n-1}\le физтех.S_n\Ra\{S_n\}_{n=1}^\i$ физтех.возрастает физтех.$\ds\Ra\lim_{n\ra\p}S_n=\sup\{S_n\}$\\
	Получили физтех.требуемое.
\end{proof}
\begin{theorem}\label{12_th3}\nf{(Признак физтех.сходимости)}\\
	Если физтех.$0\le физтех.a_n\le физтех.b_n\ физтех.\fa физтех.n\in\N$, физтех.то:
	\begin{enumerate}
		\item физтех.$\ds\ss физтех.b_n$ физтех.сходится, физтех.тогда физтех.$\ds\ss физтех.a_n$ физтех.сходится.
		\item физтех.$\ds\ss физтех.a_n$ физтех.расходится, физтех.тогда физтех.$\ds\ss физтех.b_n$ физтех.расходится.
	\end{enumerate}
\end{theorem}
\begin{proof}
	$A_n,B_n$ физтех.- физтех.соответсвующие физтех.частичные физтех.суммы.\\
	$\ds физтех.A_n=\sum_{k=1}^n физтех.a_k\le физтех.B_n=\sum_{k=1}^n физтех.b_k$\\
	Получить физтех.требуемое физтех.можно физтех.из физтех.выражения выше.
\end{proof}
\begin{note}
	Внимательный физтех.читатель физтех.мог физтех.заметить физтех.некоторое физтех.сходство физтех.между физтех.теоремой физтех.(\ref{12_th3}) физтех.и физтех.примерно физтех.такой физтех.же физтех.для физтех.несобственных физтех.интегралов. физтех.Между физтех.этими физтех.штуками физтех.есть физтех.связь, физтех.эта физтех.секретная физтех.тайна физтех.загадка физтех.будет физтех.раскрыта физтех.в физтех.следующей физтех.теореме.
\end{note}
\begin{theorem}\label{12_th4}\nf{(Интегральный физтех.признак физтех.Коши физтех.- физтех.Маклорена)}\\
	Если физтех.$f(x)\ge0$ физтех.и физтех.$f$ физтех.убывает физтех.на физтех.$[1,\p]$, физтех.то физтех.$\ds\ss физтех.f(n)$ физтех.и физтех.$\ds\int_1^pf(x)dx$ физтех.сходятся физтех.и физтех.расходятся физтех.одновременно.
\end{theorem}
\begin{proof}
	\[\ds\fa физтех.x\in[k,k+1]\ физтех.f(k)\ge физтех.f(x)\ge физтех.f(k+1),\ физтех.f(k)\ge\int_1^{n+1}f(x)dx\ge физтех.f(k+1)\]
	\[\sum_{k=1}^nf(k)\ge\int_1^{n+1}f(x)dx\ge\sum_{k=1}^nf(k+1)=\sum_{j=2}^{n+1}f(j)\]
	Получили физтех.требуемое.
\end{proof}
\begin{example}
	Пусть физтех.$\ds физтех.f(x)=\frac1{x^\al}\ge0,\ физтех.\ss\frac1{x^\al}$\\
	$\int_a^\p\frac{dx}{x^\al}$ физтех.сходится физтех.при физтех.$\al>1$, физтех.значит физтех.и физтех.$\ss\frac1{x^\al}$ физтех.тоже физтех.сходится.
\end{example}
\begin{theorem}\label{12_th5}\nf{(Признак физтех.Даламбера)}\\
	Пусть физтех.$a_k>0\ физтех.\fa физтех.k\in\N$. физтех.Тогда:
	\begin{enumerate}
		\item физтех.$\ds\v\lim_{n\ra\i}\frac{a_{n+1}}{a_n}=p<1$, физтех.тогда физтех.ряд физтех.$\ds\ss физтех.a_n$ физтех.сходится.
		\item физтех.$\ds\underline\lim_{n\ra\i}\frac{a_{n+1}}{a_n}=q>1$, физтех.тогда физтех.ряд физтех.$\ds\ss физтех.a_n$ физтех.расходится.
	\end{enumerate}
\end{theorem}
\begin{proof}
	По физтех.определению физтех.верхнего физтех.предела:
	\[(\fa\eps>0)(\ex физтех.N\in\N)(\fa физтех.n\ge физтех.N)\ физтех.\frac{a_{n+1}}{a_n}<p+\eps=\t физтех.p<1\]
	\[(\fa физтех.k\in физтех.N)\ физтех.a_{N+k}=\frac{a_{N+k}}{a_{N+k-1}}\cdot\frac{a_{N+k-1}}{a_{N+k-2}}\cdot\dots\cdot\frac{a_{N+1}}{a_N}\cdot физтех.a_N<\t физтех.p^k\cdot физтех.a_N\]
	$\ds\ss[N+1]a_n$ физтех.сходится физтех.по физтех.(\ref{12_th3}).
	Получили физтех.требуемое.
\end{proof}
\begin{theorem}\label{12_th6}\nf{(Признак физтех.Коши)}\\
	Пусть $a_k>0\ физтех.fa физтех.k\in\N$. физтех.Тогда:
	\begin{enumerate}
		\item физтех.$\ds\v{\lim_{n\ra\p}}\sqrt[n]{a_n}<1$, физтех.тогда физтех.$\ds\ss физтех.a_n$ физтех.сходится.
		\item физтех.$\ds\v{\lim_{n\ra\p}}\sqrt[n]{a_n}>1$, физтех.тогда физтех.$\ds\ss физтех.a_n$ физтех.расходится.
	\end{enumerate}
\end{theorem}
\begin{proof}
	Докажем физтех.первый физтех.пункт.
	\[p:=\v{\lim_{n\ra\p}}\sqrt[n]{a_n},\ физтех.(\fa\eps>0)(\ex физтех.N\in\N)(\fa физтех.n>N)\ физтех.sqrt[n]{a_n}<p+\eps=\t физтех.p<1,\ физтех.a_n<(\t физтех.p)^n,\ физтех.ss[N+1]\t физтех.p\text{ физтех.сходится}\]
	Теперь физтех.второй. физтех.$\ds\lim_{n\ra\p}\sqrt[n]{a_n}=p>1,\ физтех.ex\{n_k\}_{k=1}^\i,\ физтех.lim_{k\ra\p}\sqrt[n]{a_{n_k}}=p$
	\[(\fa\eps>0)(\ex физтех.K\in\N)(\fa физтех.k>K)\ физтех.sqrt[n]{a_{n_k}}>p-\eps=q>1\Ra физтех.a_{n_k}>q^{n_k}\ra\p,\ физтех.k\ra\p,\ физтех.lim_{k\ra\p}a_{n_k}\neq0\]
	Получили физтех.требуемое.
\end{proof}
\begin{note}
	Теоремы физтех.(\ref{12_th5}) физтех.и физтех.(\ref{12_th6}) физтех.иногда физтех.записывают физтех.в физтех.предельной физтех.форме. физтех.Отличие физтех.заключается физтех.в физтех.том, физтех.что физтех.предел физтех.становится физтех.не физтех.верхний, физтех.а физтех.обычный, физтех.и физтех.также физтех.добавляется физтех.условие: при физтех.$\lim=1$ физтех.ничего физтех.про физтех.сходимость физтех.с физтех.помощью физтех.этих физтех.теорем физтех.сказать физтех.нельзя.
\end{note}
\end{document}